\documentclass[11pt]{article}
\usepackage[margin=1in]{geometry}
\usepackage{amsmath,amssymb,amsthm}
\usepackage[hidelinks]{hyperref}
\newtheorem{theorem}{Theorem}
\newtheorem{proposition}[theorem]{Proposition}
\newtheorem{corollary}[theorem]{Corollary}
\newcommand{\mr}{\operatorname{mr}}
\newcommand{\Aut}{\operatorname{Aut}}
\newcommand{\F}{\mathbb F}
\title{Conjugacy-restricted quotient capacities}
\author{Internal research working proof}
\date{17 November 2026}
\begin{document}
\maketitle

Eight distinct quotient kernels onto $GL_3(\F_2)$, with meridians
mapping to involutions, suffice to prove that a knot has meridional
rank at least four.  The unrestricted two-generator quotient-kernel
capacity bound for the same finite group is 57.  The smaller threshold
uses the meridian conjugacy class rather than the ordinary group rank.
It also applies to quotient kernels of the double branched cover when
the deck involution descends to an inner automorphism.

This is a finite-group capacity criterion, without a new named knot
equality or a priority claim.  The exact enumeration below is internally
replayed; the uniform implications have written proofs and have not
received formal verification or external specialist review.

\section{Quotient kernels and restricted generating tuples}

Let $S$ be a finite group and $C\subset S$ an $\Aut(S)$-invariant set.
Let $\mathcal T_d(S,C)$ denote the ordered $d$-tuples in $C$ which
generate $S$.

\begin{proposition}\label{prop:capacity}
Suppose $\Gamma$ is generated by $d$ elements $g_1,\ldots,g_d$.
The number of distinct kernels of epimorphisms $f:\Gamma\to S$
satisfying $f(g_i)\in C$ for every $i$ is at most
\[
                    |\mathcal T_d(S,C)/\Aut(S)|.
\]
Any subgroup $A\le\Aut(S)$ gives the valid, possibly weaker, upper
bound $|\mathcal T_d(S,C)/A|$.
\end{proposition}
\begin{proof}
Evaluation on the specified generators determines a homomorphism.
Two epimorphisms $f,f'$ have the same kernel if and only if there is
$\alpha\in\Aut(S)$ such that $f'=\alpha\circ f$: both factor through
isomorphisms $\Gamma/\ker f\cong S$.  Thus distinct kernels give
distinct automorphism orbits of generating tuples.  Passing to a
subgroup can only increase the number of orbits.
\end{proof}

This elementary orbit-count argument is closely related to the
classical generating-tuple capacities for direct powers, recorded by
Burness--Crestani with credit to Hall \cite[p.~2]{bc}.  It does not
require simplicity of $S$, and the proof here does not import the
direct-power theorem.

\section{The binary three-dimensional group}

Put $S=GL_3(\F_2)$ and let $C$ be its nonidentity involutions.
Every element of $C$ is a transvection.  Indeed, if $g=I+N$ is an
involution in characteristic two, then $N^2=0$, so
$\operatorname{im}N\subset\ker N$ and $\operatorname{rank}N\le1$
in dimension three.  Nonidentity gives rank exactly one.

An invertible binary matrix is faithfully represented by its
permutation of the seven nonzero column vectors.  Enumerating all
three-column choices gives 168 matrices and 21 nonidentity involutions.
The 168 inner automorphisms, together with their compositions with
\[
                        \theta(g)=(g^{-1})^{\mathsf T},
\]
give 336 distinct automorphisms.  The centre of $S$ is trivial, as is
also checked directly.  No assertion that these exhaust $\Aut(S)$ is
needed for the upper bound.

These automorphisms are transitive on $C$.  Fix $a\in C$.
Its stabilizer among the displayed automorphisms has order 16.
Exhausting all $21^2=441$ pairs $(b,c)\in C^2$ gives
\[
\begin{array}{c|rrrrrr}
|\langle a,b,c\rangle|&2&4&6&8&24&168\\\hline
\#(b,c)&1&16&32&72&208&112.
\end{array}
\]
An automorphism fixing a generating tuple fixes every element of $S$,
so the stabilizer of any generating tuple is trivial.  The 112
generating pairs in this normalized enumeration therefore have
$112/16=7$ orbits.  Equivalently there are $21\cdot112=2352$
generating involution triples and $2352/336=7$ orbits under the
displayed automorphisms.

\begin{theorem}\label{thm:seven}
If a knot group admits eight epimorphisms onto $GL_3(\F_2)$ with
distinct kernels and with meridians mapping to involutions, then
\[
                              \mr(K)\ge4.
\]
\end{theorem}
\begin{proof}
If three meridians generate the knot group, their images under each
epimorphism form a generating involution triple.  All images are
nonidentity: a knot group is normally generated by a meridian, so
killing a meridian makes the image trivial.  Proposition~\ref{prop:capacity}
and the seven-orbit computation bound the number of distinct kernels
by seven.  A generating family of size less than three can be padded
by repeated meridians, so it is also excluded.  Meridian inverses cause
no change because their images are involutions.
\end{proof}

For comparison, the same finite enumeration finds 19152 ordered
generating pairs of arbitrary elements of $S$.  The 336 explicit
automorphisms act freely, giving the valid ordinary two-generator
capacity upper bound $19152/336=57$.  Thus the eight-kernel gate uses
additional meridional information.  It does not follow merely from
counting eight unrestricted quotient kernels of a double-cover group.

\section{When double-cover kernels supply the required quotients}

Write $G=\pi_1(S^3\setminus K)$, choose a meridian $m$, and let
\[
 \Pi=G/\langle\!\langle m^2\rangle\!\rangle,
 \qquad H=\ker(\Pi\to C_2)=\pi_1(\Sigma_2K).
\]
The map to $C_2$ sends the meridian to its nontrivial element.
The image of $m$ has order two and splits this extension, so
$\Pi=H\rtimes\langle m\rangle$.  Let $\tau$ be the automorphism of
$H$ induced by conjugation by $m$.

\begin{corollary}\label{cor:deck}
Suppose there are eight epimorphisms $f_i:H\to S=GL_3(\F_2)$
with distinct kernels.  For each $i$, require that $\ker f_i$ is
$\tau$-invariant and that the induced automorphism of $S$ is inner.
Then $\mr(K)\ge4$.
\end{corollary}
\begin{proof}
For each $f=f_i$, there is $a\in S$ with
$f(\tau h)=af(h)a^{-1}$.  Since $\tau^2=1$ and $f$ is onto,
$a^2$ is central.  The centre of $S$ is trivial, so $a^2=1$.
Consequently the assignments
\[
                 h\longmapsto f(h),\qquad m\longmapsto a
\]
define an epimorphism $\Pi\to S$.  The element $a$ cannot be the
identity, because $m$ normally generates $\Pi$ and the image is
nontrivial.  Composing with $G\to\Pi$ gives the knot-group
epimorphism required by Theorem~\ref{thm:seven}.

The kernels of these knot-group maps remain distinct: if two were
equal, their corresponding kernels in $\Pi$ would be equal, and
intersection with $H$ would give the same $\ker f_i$.  This
contradicts the hypothesis.
\end{proof}

The kernel-invariance and inner-action hypotheses are substantive.
An outer deck action does not supply this particular epimorphism
$G\to S$.  Quotients on different knot groups cannot be pooled.

\section{Replay and limits}

The standard-library checker \texttt{verify\_involution\_capacity.py}
constructs the multiplication table from literal seven-point
permutations, checks associativity, inverses and the trivial centre,
and verifies the inverse-transpose automorphism explicitly.
It checks 336 distinct automorphisms, their involution-class action,
all 441 normalized triples, full orbit marking, and all $168^2$
ordered pairs.  No historical knot-panel computation is imported.
Run
\begin{verbatim}
python3 verify_involution_capacity.py \
  --output involution_capacity_verification.json
\end{verbatim}
from this directory to record the result with a checker SHA256.
The finite enumeration supports exactly the displayed group counts.
The kernel implications are the written proofs above.  This note
does not claim a new sharp knot, completeness of a knot census,
formal verification, external review, or global mathematical priority.

\begin{thebibliography}{9}
\bibitem{bc} T. C. Burness and E. Crestani,
On the generating graph of direct powers of a simple group,
J. Algebraic Combinatorics 38 (2013), 329--350.
Author manuscript p.~2.
\url{https://seis.bristol.ac.uk/~tb13602/docs/BC_1.pdf}.
\end{thebibliography}
\end{document}
